
\section{Contribution of the Study}
The study contributes a controlled, like-for-like analytical comparison of one NBC 205:2024 Regular reference configuration and six selected plan-irregular configurations detailed using the same ready-to-use basis and evaluated under a common NBC 105:2025 demand framework. The contribution is not a proposal for automatic code extension. Rather, it is the combined use of force demand, complete drift profiles, CM--CR separation, nonlinear pushover capacity, performance-point response, final-step hinge evidence, and a study-defined Geometric Irregularity Index (GII) to examine where the response of selected irregular plans departs from the Regular reference.

A second contribution is the explicit separation between geometric description and seismic performance. Configurations with the same area-deficit GII are retained as distinct models, allowing the study to test whether equal GII values reproduce equal response. A third contribution is the evidence-based comparison of a representative Draft Final NBC 205:2012 building, the NBC 205:2024 Regular benchmark, and a separate four-storey height-extension case without mixing the different types of available analysis evidence. The resulting conclusions therefore identify practical applicability boundaries and evidence gaps rather than claiming universal approval of RUD detailing for irregular buildings.
